\subsubsection*{Solution}

\begingroup
\color{blue}

The derivation below applies to the minimal local realization specified by the
blue additions to the problem statement.

\paragraph*{Local Ward identities}

Local charge conservation reads

\[
\partial_t n+\partial_x J^{(1)}=0.
\]

For a localized fluctuation, define

\[
\delta N=\int\mathrm{d}x\,n,
\qquad
\delta D=\int\mathrm{d}x\,x n,
\qquad
\delta Q=\int\mathrm{d}x\,x^2 n.
\]

All boundary terms vanish. Integrating the continuity equation by parts gives

\[
\frac{\mathrm{d}\delta D}{\mathrm{d}t}
=\int\mathrm{d}x\,J^{(1)}.
\]

To implement dipole conservation locally without introducing an independent
dipole density, write

\[
J^{(1)}=\partial_x J^{(2)}.
\]

The quadrupole moment then satisfies

\[
\frac{\mathrm{d}\delta Q}{\mathrm{d}t}
=2\int\mathrm{d}x\,x J^{(1)}
=-2\int\mathrm{d}x\,J^{(2)}.
\]

Quadrupole conservation requires the last integral to vanish. With no
independent quadrupole density, write

\[
J^{(2)}=\partial_x J^{(3)}.
\]

Thus the three moment constraints are represented locally by

\[
\partial_t n+\partial_x^3 J^{(3)}=0.
\tag{1}
\]

This is the minimal Ward hierarchy. Keeping intrinsic dipole or quadrupole
densities would enlarge the hydrodynamic theory and can introduce additional
modes.

\paragraph*{Free energy and constitutive relations}

To quadratic order in the fields and at the lowest spatial derivative order,
the free energy is

\[
F[n,\theta]
=\frac12\int\mathrm{d}x\,
\left[
\frac{n^2}{\chi}
+\kappa(\partial_x\theta)^2
\right].
\]

Writing $f$ for the free-energy density, the chemical potential and
stiffness current are

\[
\mu=\frac{\delta F}{\delta n}=\frac{n}{\chi},
\qquad
J^{(3)}
=\frac{\partial f}{\partial(\partial_x\theta)}
=\kappa\partial_x\theta.
\]

Using the Ward hierarchy,

\[
J^{(2)}=\partial_x J^{(3)}=\kappa\partial_x^2\theta,
\qquad
\frac{\delta F}{\delta\theta}=-J^{(2)}.
\]

Equation (1) is equivalently
$\partial_t n=-\partial_x^2J^{(2)}$. The free energy therefore changes at the
rate

\[
\begin{aligned}
\frac{\mathrm{d}F}{\mathrm{d}t}
&=\int\mathrm{d}x\,
\left(
\mu\,\partial_t n
+\frac{\delta F}{\delta\theta}\partial_t\theta
\right)\\
&=-\int\mathrm{d}x\,J^{(2)}
\left(\partial_t\theta+\partial_x^2\mu\right).
\end{aligned}
\]

The factor in parentheses is the thermodynamic force conjugate to the
rank-two current. With a single local scalar Onsager channel, the leading
dissipative constitutive relation is

\[
J^{(2)}
=\sigma\left(\partial_t\theta+\partial_x^2\mu\right),
\qquad \sigma>0.
\]

It gives

\[
\frac{\mathrm{d}F}{\mathrm{d}t}
=-\sigma\int\mathrm{d}x\,
\left(\partial_t\theta+\partial_x^2\mu\right)^2
\leq0.
\]

Combining these relations with $n=\chi\mu$ gives

\[
\chi\partial_t\mu+\kappa\partial_x^4\theta=0,
\tag{2}
\]

\[
\sigma\partial_t\theta
+\sigma\partial_x^2\mu
-\kappa\partial_x^2\theta=0.
\tag{3}
\]

\paragraph*{Schwinger-Keldysh retarded kernel}

The retarded poles follow from Eqs.~(2) and (3). The corresponding $ar$ part
of the quadratic Schwinger-Keldysh action can be written as

\[
\begin{aligned}
S_{ar}^{(2)}=\int\mathrm{d}t\,\mathrm{d}x\,\bigl[&
\mu_a\left(\chi\partial_t\mu_r
+\kappa\partial_x^4\theta_r\right)\\
&+\theta_a\left(
\sigma\partial_t\theta_r
+\sigma\partial_x^2\mu_r
-\kappa\partial_x^2\theta_r
\right)\bigr].
\end{aligned}
\]

Here $r$ denotes the physical, or average, field and $a$ the response, or
difference, field. The $aa$ terms fixed by KMS encode noise but do not enter
the tree-level retarded kernel. Using
$\exp(-i\omega t+ikx)$, Eqs.~(2) and (3) become

\[
K_R(\omega,k)
\begin{pmatrix}
\mu\\
\theta
\end{pmatrix}
=0,
\qquad
K_R(\omega,k)=
\begin{pmatrix}
-i\chi\omega & \kappa k^4\\
-\sigma k^2 & -i\sigma\omega+\kappa k^2
\end{pmatrix}.
\]

The pole condition is

\[
\det K_R
=-\chi\sigma\omega^2
-i\chi\kappa k^2\omega
+\kappa\sigma k^6
=0,
\]

or

\[
\omega^2
+i\frac{\kappa}{\sigma}k^2\omega
-\frac{\kappa}{\chi}k^6=0.
\tag{4}
\]

Its two roots are

\[
\omega_\pm(k)
=-\frac{i\kappa k^2}{2\sigma}
\pm\sqrt{
\frac{\kappa k^6}{\chi}
-\frac{\kappa^2k^4}{4\sigma^2}
}.
\tag{5}
\]

\paragraph*{Long-wavelength behavior}

At fixed finite $\chi$, $\kappa$, and $\sigma$, Eq.~(5) has the expansions

\[
\omega_{\mathrm{slow}}
=-i\frac{\sigma}{\chi}k^4
-i\frac{\sigma^3}{\kappa\chi^2}k^6
+O(k^8),
\]

\[
\omega_{\mathrm{fast}}
=-i\frac{\kappa}{\sigma}k^2
+i\frac{\sigma}{\chi}k^4
+i\frac{\sigma^3}{\kappa\chi^2}k^6
+O(k^8).
\]

Both poles lie in the lower half-plane for positive coefficients. The slow
$k^4$ mode is density dominated, whereas the $k^2$ mode describes relaxation
of the quadrupole Goldstone current. The limit $\sigma\to0$ is outside this
truncation because the omitted Goldstone inertia then becomes important.
Equation~(5) is exact within the quadratic EFT; only its small-$k$ behavior is
controlled by hydrodynamics.

\endgroup
